\apendice{Especificación de Requisitos}

\section{Explicación casos de uso.}

Un caso de uso es un concepto utilizado en el desarrollo de \textit{software} para describir una secuencia de acciones a realizar para lograr un resultado observable para el usuario \cite{Manuel-Cillero}. A continuación, se especifican los casos de uso y se muestran los diagramas de secuencias de la aplicación UV-Dermis.

Para el acceso y visualización del cotenido de la pestaña de mapas de variabes se es sigue el proceso detallado en el diagrama \ref{fig:interna3}. De la misma manera, para el acceso de riesgo acumulado esta la figura \ref{fig:interna2} y para el acceso al recomendador  la figura \ref{fig:interna1}.

\begin{figure}[H]
        \centering
        \includegraphics[width=0.9\textwidth]{img/riesgoacum.png}
        \caption{Diagrama de secuencias, acceso al cálculo de riesgo acumulado.}
        \label{fig:interna2}
\end{figure}

\begin{figure}[H]
        \centering
        \includegraphics[width=0.9\textwidth]{img/mapas.png}
        \caption{Diagrama de secuencias, acceso a mapa de variables meteorológicas.}
        \label{fig:interna3}
\end{figure} 

\begin{figure}[H]
        \centering
        \includegraphics[width=0.9\textwidth]{img/recomen.png}
        \caption{Diagrama de secuencias, acceso a la recomendación.}
        \label{fig:interna1}
\end{figure}

En la figura \ref{fig:interna}, asociada al caso de uso \ref{fig:cu5}, se observa que la aplicación para mostrar los datos meteorológicos no vas mas allá del entorno local.
 \begin{figure}[H]
        \centering
        \includegraphics[width=0.9\textwidth]{img/datospesta.png}
        \caption{Diagrama de secuencias, visualización de datos.}
        \label{fig:interna}
    \end{figure}



\begin{table}[p]
	\centering
	\begin{tabularx}{\linewidth}{ p{0.21\columnwidth} p{0.71\columnwidth} }
		\toprule
		\textbf{CU-1}    & \textbf{Acceso contenido informativo}\\
		\toprule
		\textbf{Versión}              & 1.0    \\
		\textbf{Autor}                & Nisrine Fariss Lamine \\
		
		\textbf{Descripción}          & Información relevante para entender el contexto de la aplicación \\
		\textbf{Precondición}         & Acceder a las pestañas del contenido infromativo y marco teórico de la calculadora \\
		\textbf{Acciones}             &
		\begin{enumerate}
			\def\labelenumi{\arabic{enumi}.}
			\tightlist
			\item En el módulo del contenido informativo acceder a las diferentes subpestañas.
			\item En la pestaña de riesgo acumulado acceder al marco teórico de la calculadora.
		\end{enumerate}\\
		\textbf{Postcondición}        & El usuario obtiene información del contexto y fundamentos científicos. \\
	
		\textbf{Importancia}          & Media \\
		\bottomrule
	\end{tabularx}
	\caption{CU-1 Acceso contenido informativo.}
\end{table}

\begin{table}[p]
	\centering
	\begin{tabularx}{\linewidth}{ p{0.21\columnwidth} p{0.71\columnwidth} }
		\toprule
		\textbf{CU-2}    & \textbf{Visualización de mapas}\\
		\toprule
		\textbf{Versión}              & 1.0    \\
		\textbf{Autor}                & Nisrine Fariss Lamine \\
		
		\textbf{Descripción}          & Exploración la variabilidad geográfica de las variables. \\
		\textbf{Precondición}         & Acceso al módulo de mapa de variables. Conexión a Internet. \\
		\textbf{Acciones}             &
		\begin{enumerate}
			\def\labelenumi{\arabic{enumi}.}
			\tightlist
			\item Visitar las subpestañas del módulo.
			
		\end{enumerate}\\
		\textbf{Postcondición}        & Se muestran los mapas \\
		\textbf{Excepciones}          & Si la AEMET/INE no responden se muestra una ventana gris. \\
		\textbf{Importancia}          & Alta \\
		\bottomrule
	\end{tabularx}
	\caption{CU-2 Visualización de mapas.}
\end{table}


\begin{table}[p]
	\centering
	\begin{tabularx}{\linewidth}{ p{0.21\columnwidth} p{0.71\columnwidth} }
		\toprule
		\textbf{CU-1}    & \textbf{Cálculo riesgo acumulado}\\
		\toprule
		\textbf{Versión}              & 1.0    \\
		\textbf{Autor}                & Nisrine Fariss Lamine \\
		
		\textbf{Descripción}          & Estimación del riesgo acumulado de desarrollar melanoma para el usuario \\
		\textbf{Precondición}         & Acceso a la pestaña de riesgo acumulado. Disponibilidad de datos. \\
		\textbf{Acciones}             &
		\begin{enumerate}
			\def\labelenumi{\arabic{enumi}.}
			\tightlist
			\item Introducir los parámetros del usuario.
			\item Ejecución del cálculo.
		\end{enumerate}\\
		\textbf{Postcondición}        & Disposición de una estimación orientativa sobre los factores de riesgo actuales. \\
		\textbf{Excepciones}          & Ausencia de datos. \\
		\textbf{Importancia}          & Alta  \\
		\bottomrule
	\end{tabularx}
	\caption{CU-3 Cálculo riesgo acumulado.}
\end{table}



\begin{table}[p]
	\centering
	\begin{tabularx}{\linewidth}{ p{0.21\columnwidth} p{0.71\columnwidth} }
		\toprule
		\textbf{CU-4}    & \textbf{Obtener recomendación}\\
		\toprule
		\textbf{Versión}              & 1.0    \\
		\textbf{Autor}                & Nisrine Fariss Lamine \\
		
		\textbf{Descripción}          & Ofrece al usuario medidas de prevención según su fototipo de piel. \\
		\textbf{Precondición}         & Acceder a la pestaña del recomendador \\
		\textbf{Acciones}             &
		\begin{enumerate}
			\def\labelenumi{\arabic{enumi}.}
			\tightlist
			\item Elección de provincia y fototipo de piel.
			\item Obtener recomendación
		\end{enumerate}\\
		\textbf{Postcondición}        & Se obtiene una recomendación preventiva e información sobre la temperatura e índice UV de su provincia. \\
		\textbf{Excepciones}          & Ausencia de datos. \\
		\textbf{Importancia}          & Alta \\
		\bottomrule
	\end{tabularx}
	\caption{CU-4 Obtener recomendación.}
\end{table}



\begin{table}[H]
	\centering
	\begin{tabularx}{\linewidth}{ p{0.21\columnwidth} p{0.71\columnwidth} }
		\toprule
		\textbf{CU-5}    & \textbf{Exploración de historico de datos meteorológicos}\\
		\toprule
		\textbf{Versión}              & 1.0    \\
		\textbf{Autor}                & Nisrine Fariss Lamine \\
		
		\textbf{Descripción}          & Facilita al usuario visualización de la evolución temporal de temperaturas e índices UV de las provincias, ambas seleccionables. \\
		\textbf{Precondición}         & Acceso a la pestaña de datos meteorológicos. Existencia del conjunto de datos pasados. \\
		\textbf{Acciones}             &
		\begin{enumerate}
			\def\labelenumi{\arabic{enumi}.}
			\tightlist
			\item Selección de provincia/s y variable/s.
		\end{enumerate}\\
		\textbf{Postcondición}        & Gráfico evolutivo de las variables. Posibilidad de descarga \\
		\textbf{Excepciones}          &  \\
		\textbf{Importancia}          & Media \\
		\bottomrule
	\end{tabularx}
	\caption{CU-5 Exploración de historico de datos meteorológicos.}
    \label{fig:cu5}

    \end{table}


   